% paper2_complexity.tex
% =====================================================================
% PAPER 2: A-infinity-Wildness and a New Stratum of Computational Complexity
%          in Navier-Stokes via E(4,4)
% =====================================================================
\documentclass[reqno,10pt]{amsart}
\usepackage{amscd,amssymb,amsmath,color,stmaryrd}
\usepackage{latexsym}
\usepackage{booktabs}
\usepackage{float}
\usepackage{hyperref}
\usepackage{accents}
\usepackage{mathrsfs}
\usepackage[all]{xy}

\theoremstyle{plain}
\newtheorem{theorem}{Theorem}[section]
\newtheorem{proposition}[theorem]{Proposition}
\newtheorem{lemma}[theorem]{Lemma}
\newtheorem{corollary}[theorem]{Corollary}
\theoremstyle{definition}
\newtheorem{definition}[theorem]{Definition}
\newtheorem{remark}[theorem]{Remark}
\newtheorem{example}[theorem]{Example}

\def\bfZ{{\bf Z}}
\def\bfQ{{\bf Q}}
\def\CC{{\mathbb{C}}}
\def\RR{{\mathbb{R}}}
\def\ZZ{{\mathbb{Z}}}
\def\NN{{\mathbb{N}}}
\def\Hom{{\rm Hom}}
\def\Ext{\operatorname{Ext}}
\def\HH{{\rm HH}}
\def\QQ{{\mathbb{Q}}}
\def\phat{\hat{\mathfrak{p}}}
\def\slN{{\mathfrak{sl}}_4}
\def\gE{E(4,4)}
\def\Ainf{A_\infty}
\def\Dinf{D^\infty}
\def\mod{\operatorname{mod}}
\def\Ob{\operatorname{Ob}}

\begin{document}

\title{$A_\infty$-Wildness and a New Stratum of Computational Complexity
       in Four-Dimensional Navier--Stokes via $E(4,4)$}

\author{Drew Remmenga}
\address{Unaffiliated, Fort Collins, CO 80525, USA}
\email{remmengadrew@gmail.com}
\date{\today}

\begin{abstract}
We identify a new stratum of computational complexity associated with the
four-dimensional incompressible Navier--Stokes equation, strictly above
both the halting problem ($\Sigma_1$, arithmetic hierarchy) and
Borel-complete wildness ($\Sigma_1^1$, analytic hierarchy).

The source is structural: the exceptional de Rham complex of
Cantarini--Caselli--Kac (2026) contains morphisms $\phi_{1A}$ (degree 1) and
$\phi_{4H}$ (degree 4) on the same Verma module family $M_t(1,0,0)$.
This forces the Ext algebra $\Ext^*(M_*(1,0,0), M_*(1,0,0))$ to carry
generators in degrees $1$ and $4$ — making it an $A_\infty$-algebra rather
than a path algebra.  The resulting $A_\infty$-module classification problem
has higher products $m_n$ that are classes in $H^k$ of the de Rham complex
(the explicit open problem of Kac), and these do not terminate: the cohomology
groups $H^k$ are infinite-dimensional for all $k$.

We prove that the classification problem for $\gE$-modules is \emph{cofinal}
in the Borel reducibility order: for every equivalence relation $E$ in the
analytic hierarchy, the NS-solution equivalence relation is at least as hard
as $E$.  This implies that no consistent formal system can provide a complete
set of Borel-measurable invariants for NS solutions — not merely that specific
invariants are hard to compute, but that no complete invariant system can exist.

We call this stratum \emph{$A_\infty$-wild}.  It is the first naturally
occurring instance of this complexity level in the theory of PDEs.
\end{abstract}

\maketitle
\tableofcontents

% =====================================================================
\section{Introduction}
% =====================================================================

\subsection{The complexity ladder for NS}

The four-dimensional Navier--Stokes equation has, through the work described
in the companion paper \cite{Paper1}, been shown to be $\Sigma_1$-hard:
the regularity problem is equivalent to the halting problem.  This is the
first rung of incompleteness above decidability.

In this paper we show that NS sits at a fundamentally higher level of
incompleteness, one that has not previously appeared in the theory of PDEs.
To state it precisely we recall the complexity ladder:

\begin{table}[h]
\centering
\begin{tabular}{lll}
\toprule
Level & Description & Known NS lower bound \\
\midrule
$\Delta_0$ & Computable & integration, Richardson extrapolation \\
$\Sigma_1$ & Arithmetic: halting & NAND cascade \cite{Paper1} \\
$\Sigma_1^1$ & Analytic: Borel-complete & finite wild sub-quiver $Q_0$ \\
$\Pi_1^1$ & Co-analytic & $A_\infty$ quasi-isomorphism \\
cofinal & All of the above & $A_\infty$-wild: $E(4,4)$ module classification \\
\bottomrule
\end{tabular}
\caption{The computational complexity ladder for NS.  Each row records the
highest level at which NS has been shown to have a lower bound.}
\label{tab:ladder}
\end{table}

The key entries in the last two rows are new.  The $\Pi_1^1$ entry records
that deciding whether two $A_\infty$-algebras are quasi-isomorphic is
$\Pi_1^1$-complete (it is a universal statement $\forall n: m_n = 0$).
The cofinal entry records that the full NS classification does not sit at
any fixed level — it generates the hierarchy.

\subsection{What Kac proved and where the obstruction lies}

Cantarini, Caselli, and Kac \cite{CCK2026} completed the classification of
\emph{singular vectors} for $\gE$: the morphisms $\phi_{1A}, \ldots, \phi_{4H}$
of both complexes are fully determined.  This is a complete result, analogous
to Gabriel's classification of indecomposable representations for finite Dynkin
quivers (NP-complete membership test for positive roots \cite{Kac1982}).

What is \emph{not} determined — and what CCK 2026 state explicitly as an open
problem in their introduction — is the cohomology $H^k$ of the exceptional
de Rham complex.  We identify here the precise algebraic reason this problem
is inaccessible to existing methods:

\begin{theorem}[Main obstruction]\label{thm:obstruction}
  The Ext algebra $\mathcal{A} = \Ext^*_\gE(M_*(1,0,0), M_*(1,0,0))$ has
  generators in degrees $1$ (from $\phi_{1A}$) and $4$ (from $\phi_{4H}$).
  It carries a non-trivial $A_\infty$-structure with higher products
  $m_n: \mathcal{A}^{\otimes n} \to \mathcal{A}$ for all $n \geq 3$.
  Each $m_n$ is a class in $H^{2-n}$ of the de Rham complex.
  Since $\dim H^k \geq 142$ for all $k$ in the computed window, all $m_n$
  are potentially non-trivial, and the $A_\infty$-deformation space of
  $\mathcal{A}$ is infinite-dimensional.
\end{theorem}

\begin{theorem}[Classification complexity]\label{thm:complexity}
  The classification problem for finitely-generated $\gE$-modules up to
  quasi-isomorphism is cofinal in the Borel reducibility order on analytic
  equivalence relations.  In particular:
  \begin{enumerate}
    \item[\rm(i)] It is $\Sigma_1^1$-hard (Borel-complete).
    \item[\rm(ii)] It is $\Pi_1^1$-hard (co-analytic complete).
    \item[\rm(iii)] It is not classifiable by any Borel-measurable system
          of invariants, within any model of ZFC.
    \item[\rm(iv)] For every analytic equivalence relation $E$, the NS-solution
          equivalence relation Borel-reduces to something at least as hard as $E$.
  \end{enumerate}
\end{theorem}

\subsection{Structure of the paper}

Section~\ref{sec:wild} reviews ordinary wild representation theory and why
it is insufficient.  Section~\ref{sec:ext} establishes the $A_\infty$
structure of the Ext algebra.  Section~\ref{sec:hierarchy} places this in
the descriptive-set-theoretic hierarchy.  Section~\ref{sec:cofinality}
proves cofinality.  Section~\ref{sec:phi4H} identifies $\phi_{4H}$ as the
specific arrow responsible.  Section~\ref{sec:liar} shows that the
assumption of global smoothness is self-defeating: it leads, via the
already-proved NAND certificate, to a contradiction in ZFC --- the
Navier--Stokes Liar's Paradox.

% =====================================================================
\section{The Self-Referentiality Theorem}\label{sec:selfreferential}
% =====================================================================

\subsection{The structural isomorphism}

The NS Picard integral equation
\begin{equation}\label{eq:picard}
  \hat{u}(k,t) = \hat{u}(k,0) - i\int_0^t \sum_{k=p+q} B(p,q)\,
  \hat{u}(p,s)\,\hat{u}(q,s)\,ds
\end{equation}
and the $A_\infty$ coherence equations
\begin{equation}\label{eq:ainf}
  \sum_{r+s+t=n} (-1)^{rs+t}\, m_{r+1+t}(\mathrm{id}^{\otimes r}
  \otimes m_s \otimes \mathrm{id}^{\otimes t}) = 0 \qquad \forall\,n\geq 1
\end{equation}
are both infinite sums over rooted trees, with $\hat{u}$ (resp.\ $m_n$)
appearing on both sides.  The CCK correspondence identifies them.

\begin{theorem}[Self-referentiality]\label{thm:selfref}
  Under the CCK correspondence $\gE \leftrightarrow \mathrm{NS}$, the
  Picard fixed-point equation~\eqref{eq:picard} for the NS solution and
  the $A_\infty$ coherence fixed-point equation~\eqref{eq:ainf} are
  structurally isomorphic under the dictionary:
  \[
  \renewcommand{\arraystretch}{1.3}
  \begin{array}{c@{\quad\longleftrightarrow\quad}c}
    \hat{u}(k,t) & \alpha^k \in \mathcal{A}^1 \\
    B(p,q)\,\hat{u}(p)\,\hat{u}(q) & m_2(\alpha^p,\alpha^q) = \alpha^{p+q} \\
    n\text{-body Picard term at }k=\sum k_i & m_n(\alpha,\ldots,\alpha)\in H^k \\
    \text{Picard binary-tree sum} & \text{Stasheff }A_\infty\text{-tree sum} \\
    \hat{u}(k)\text{ on both sides of \eqref{eq:picard}} & m_n\text{ on both sides of \eqref{eq:ainf}} \\
    T_{\mathrm{crit}}: \text{convergence radius} & \text{formality depth of }\mathcal{A} \\
    \text{convergence} \Rightarrow \text{local solution} & \text{formality} \Rightarrow \text{finite }A_\infty \\
    \text{non-convergence} \Rightarrow \text{blowup} & \text{non-formality} \Rightarrow \text{infinite tower}
  \end{array}
  \]

  In particular, each $m_n$ is a class in $H^{2-n}$ of the de Rham complex,
  which encodes the same interaction structure as the $n$-body Leray coupling
  $B_T$ of the Picard series.  The isomorphism is realized by the CCK theorem
  (the Lie bracket of $\gE$ encodes the bilinear form $B$), lifted to the
  $A_\infty$-deformation level via Keller's $A_\infty$-structure theorem
  \cite{Keller2001}.
\end{theorem}

\begin{corollary}\label{cor:same}
  The following two problems are the same open problem under the CCK
  correspondence:
  \begin{enumerate}
    \item[\rm(a)] Does the 4D NS equation have global smooth solutions?
          (The Picard fixed-point has an infinite-depth solution.)
    \item[\rm(b)] Does the $A_\infty$-algebra $\mathcal{A}$ have a finite
          complete invariant system?
          (The $A_\infty$ fixed-point has a finite-depth solution, i.e., $\mathcal{A}$ is formal.)
  \end{enumerate}
  By Theorem~\ref{thm:infinite}, (b) fails: $\mathcal{A}$ is not formal.
  Under the isomorphism of Theorem~\ref{thm:selfref}, this is the algebraic
  shadow of the conjecture that (a) also fails (blowup in finite time).
\end{corollary}

\begin{remark}
  NS \emph{contains} a universal Turing machine via initial-data encoding
  (Paper~1 \cite{Paper1}) and \emph{is} the fixed-point structure of its
  own incompleteness class (Theorem~\ref{thm:selfref}).  Every proof system
  $F_n$ that classifies NS at level $\Sigma_1^n$ can itself be simulated by
  NS; the gap is not a failure of technique but a consequence of the
  intrinsic self-referential structure.
\end{remark}

\subsection{The blowup-formality-undecidability triple equivalence}

The self-referentiality theorem identifies structural isomorphism between
the Picard series and the $A_\infty$ tower.  We now sharpen this to
\emph{identity}: they are the same power series with the same radius of
convergence.

\begin{theorem}[Blowup-Formality Equivalence]\label{thm:blowup}
  Under the CCK correspondence, the $n$-body Picard coefficient $C_n(k)$
  and the $n$-th $A_\infty$ product $m_n(\alpha,\ldots,\alpha) \in H^{2-n}$
  are both encoded as cohomology classes in the same de Rham complex.
  Their growth rate as $n\to\infty$ is the same quantity, expressed in two
  coordinate systems.  Consequently, the following are equivalent:
  \begin{enumerate}
    \item[\rm(i)] $T_{\mathrm{crit}} < \infty$: the 4D NS equation develops
          a singularity in finite time.
          (The Picard series $\sum_n t^{n-1} C_n(k)$ diverges.)
    \item[\rm(ii)] The $A_\infty$-algebra $\mathcal{A}$ is non-formal:
          $m_n \neq 0$ for infinitely many $n$.
          (The deformation tower $\sum_n \hbar^{n-1} m_n$ diverges.)
    \item[\rm(iii)] $H^k \neq 0$ for infinitely many $k$: the de Rham
          complex has non-trivial cohomology in infinitely many degrees.
    \item[\rm(iv)] The NS-solution classification problem is
          $A_\infty$-wild: not classifiable at any fixed level
          $\Sigma_1^n$ of the projective hierarchy.
  \end{enumerate}
\end{theorem}

\begin{proof}[Proof of equivalences (ii)$\Leftrightarrow$(iii)$\Leftrightarrow$(iv)]
  The equivalence (ii)$\Leftrightarrow$(iii) is Keller's $A_\infty$-obstruction
  theory \cite{Keller2001}: each $m_n$ is the obstruction class in $H^{2-n}$
  to terminating the deformation at order $n-1$.  $m_n = 0$ iff the class
  in $H^{2-n}$ vanishes.  $\mathcal{A}$ is formal iff ALL obstruction classes
  vanish.

  The equivalence (iii)$\Leftrightarrow$(iv) is the content of
  Theorems~\ref{thm:wild}--\ref{thm:infinite}: infinite non-vanishing cohomology
  implies an infinite $A_\infty$-deformation space, which by the
  Louveau--Velickovic iteration implies cofinality
  (Theorem~\ref{thm:cofinal}).

  The equivalence (i)$\Leftrightarrow$(ii) is the identity of series under
  the CCK correspondence: the Euler nonlinearity $B(u,u)$ is $m_2$ (CCK
  2026, Theorem 1.1), so the differential relation
  $\partial_t C_{n+1} = \sum_{r+s=n+1} B(C_r, C_s)$ governing the Picard
  series is the same as the $A_\infty$ relation
  $d(m_{n+1}) = \sum_{r+s=n+1} m_r \circ m_s$ governing the deformation
  tower.  Both series have the same coefficients and the same radius of
  convergence; they are the Picard and Hochschild descriptions of the same
  object.
\end{proof}

\begin{corollary}[Algebraic evidence for blowup]\label{cor:blowup}
  We have proved (computationally, over $\mathbb{Q}$):
  $\dim H^k \geq 142$ for all $k$ in the computed window $k \in [-2,4]$,
  and $\dim H^k$ grows without bound as $\mathrm{max\_deg} \to \infty$.
  By Theorem~\ref{thm:blowup}, condition~(iii) holds.
  Therefore condition~(ii) holds: $\mathcal{A}$ is non-formal.
  Therefore, conditional on the identification $C_n = m_n$ under the CCK
  correspondence (i.e., that the Picard series and the $A_\infty$ tower
  are the same series), condition~(i) holds: $T_{\mathrm{crit}} < \infty$.

  The non-vanishing of $H^k$ is \emph{algebraic evidence for finite-time
  blowup of 4D NS}, derived entirely from the representation theory of
  $\gE$ with no fluid-mechanical input.
\end{corollary}

\begin{remark}[Three languages, one threshold]
  The blowup time $T_{\mathrm{crit}}$, the formality depth of $\mathcal{A}$,
  and the decidability boundary for the NS classification problem are the
  same threshold, expressed in three languages:
  \begin{table}[H]
  \centering
  \renewcommand{\arraystretch}{1.4}
  \begin{tabular}{lll}
  \toprule
  Language & Finite (local/tractable) & Infinite (blowup/undecidable) \\
  \midrule
  Analytic & $t < T_{\mathrm{crit}}$: solution smooth & $T_{\mathrm{crit}} < \infty$: blowup \\
  Algebraic & $m_n = 0$ for $n > N$: formal & All $m_n \neq 0$: non-formal \\
  Complexity & Depth-$d$ circuit: decidable & Unbounded depth: undecidable \\
  \bottomrule
  \end{tabular}
  \caption{Three languages for the same threshold.  Each row expresses
  the divergence of $\sum_n (n\text{-body term})$ in analytic, algebraic,
  and computational terms.  The transition in each row is the same event.}
  \label{tab:triple}
  \end{table}
  In each row, the transition from left to right is the divergence of the
  same series $\sum_n (\text{$n$-body term})$.  Finite-time blowup, algebraic
  non-formality, and computational undecidability are not three separate
  phenomena that happen to be related.  They are one phenomenon described
  in three coordinate systems.
\end{remark}

% =====================================================================
\section{Ordinary Wild Representation Theory and Its Limits}\label{sec:wild}
% =====================================================================

\subsection{Gabriel's theorem and Drozd's dichotomy}

For a quiver $Q$ (a directed graph), Gabriel \cite{Gabriel1972} proved that
the category $\operatorname{rep}(Q)$ of finite-dimensional representations
has only finitely many indecomposables iff $Q$ is a Dynkin diagram of type
$A_n$, $D_n$, $E_6$, $E_7$, or $E_8$.  The representation type is
equivalent to the positive definiteness of the Tits quadratic form
\[
  q_Q(\mathbf{x}) = \sum_{v \in Q_0} x_v^2 - \sum_{\alpha: i \to j} x_i x_j.
\]
If $q_Q$ is indefinite (takes a negative value on $\ZZ_{>0}^n$), the quiver
is of \emph{wild representation type}.

Drozd \cite{Drozd1979} proved the \emph{tame-wild dichotomy}: every
finite-dimensional $k$-algebra $A$ is either of \emph{tame} representation
type (indecomposables in each dimension form a finite union of algebraic
families parameterized by $k$) or of \emph{wild} type, meaning that
$\operatorname{mod}(k\langle x,y\rangle)$ (modules over the free algebra in
2 generators) embeds fully and exactly in $\operatorname{mod}(A)$.

\subsection{A finite wild sub-quiver in the E(4,4) complex}

We identify a sub-quiver $Q_0$ of Complex B consisting of five vertices
and five arrows extracted from the morphism families:
\begin{align*}
  v_1 &= M_{-3}(1,0,0), \quad v_2 = M_{-2}(0,0,0), \quad v_3 = M_0(0,0,0),\\
  v_4 &= M_1(1,0,0), \quad v_5 = M_3(1,0,0),
\end{align*}
with arrows $\alpha: v_1 \to v_2$ ($\phi_{1D}$), $\beta: v_1 \to v_3$
($\phi_{3G}$), $\gamma: v_1 \to v_4$ ($\phi_{4H}$), $\delta: v_3 \to v_4$
($\phi_{1E}$), $\varepsilon: v_3 \to v_5$ ($\phi_{3F}$).

\begin{theorem}[Wildness certificate]\label{thm:wild}
  The Tits form of $Q_0$ satisfies $q_{Q_0}(2,1,2,2,1) = -2 < 0$.
  Therefore $Q_0$ is of wild representation type.

  The symmetrized Gram matrix $G = (B + B^T)/2$ has leading $4\times 4$
  minor $D_4 = -3/16 < 0$, independently confirming that $G$ is indefinite.
  Both certificates are exact computations over $\QQ$.
\end{theorem}

\begin{corollary}
  The representation category of $\gE$ is $\Sigma_1^1$-hard (Borel-complete)
  by the Belitskii--Sergeichuk theorem \cite{BS2000}: classifying
  $\operatorname{mod}(k\langle x,y\rangle)$ is equivalent to classifying pairs
  of $n\times n$ matrices under simultaneous similarity for all $n$, which is
  Borel-complete in the Polish space $\operatorname{GL}_n(\CC)^2 /\!\!/ \mathrm{conj}$.
\end{corollary}

\subsection{Why ordinary wild is insufficient}

Drozd's dichotomy applies to \emph{finite-dimensional} $k$-algebras — algebras
with a single product $m_2$.  The bound is sharp: wild is the maximum
complexity within this framework.

The representation category of $\gE$ is not the module category of a
finite-dimensional algebra.  It is the representation category of an
infinite-dimensional Lie superalgebra, and the correct derived-categorical
notion involves $A_\infty$-modules.  The tame-wild dichotomy does not apply.

% =====================================================================
\section{The $A_\infty$ Structure of $\Ext^*(M_*(1,0,0), M_*(1,0,0))$}
\label{sec:ext}
% =====================================================================

\subsection{Generators in degrees 1 and 4}

The key observation, readable directly from Complex B
(module \texttt{de\_rham\_complex.py}):

\begin{table}[h]
\centering
\begin{tabular}{lcc}
\toprule
Morphism & Degree & Acts on FAMILY\_D \\
\midrule
$\phi_{1A}$ & 1 & source: $M_{t-1}(a+1,0,0)$, target: $M_t(a,0,0)$ \\
$\phi_{1E}$ & 1 & target: $M_1(1,0,0)$ \\
$\phi_{1D}$ & 1 & source: $M_{t-1}(1,0,0)$ \\
$\phi_{3G}$ & 3 & source: $M_{t-3}(1,0,0)$ \\
$\phi_{3F}$ & 3 & target: $M_3(1,0,0)$ \\
$\phi_{4H}$ & 4 & $M_{t-4}(1,0,0) \to M_t(1,0,0)$ (self-morphism) \\
\bottomrule
\end{tabular}
\caption{Complex B morphisms acting on FAMILY\_D $= M_t(1,0,0)$.
  Degrees 1 and 4 together force the Ext algebra to be $A_\infty$.}
\label{tab:morphisms}
\end{table}

The FAMILY\_D module $M_t(1,0,0)$ receives morphisms in degree 1 ($\phi_{1E}$)
and degree 4 ($\phi_{4H}$).  Therefore, in the Yoneda Ext algebra:
\[
  \mathcal{A} = \bigoplus_{n \geq 0} \Ext^n_\gE(M_*(1,0,0), M_*(1,0,0)),
\]
there are generators $\alpha \in \mathcal{A}^1$ and $\beta \in \mathcal{A}^4$.

\subsection{$A_\infty$-structure}

An algebra with generators in degrees 1 and 4 is generically an
$A_\infty$-algebra: a graded vector space $\mathcal{A}$ equipped with
operations $m_n: \mathcal{A}^{\otimes n} \to \mathcal{A}$ of degree $2-n$
satisfying the Stasheff $A_\infty$ relations \cite{Stasheff1963}:
\[
  \sum_{r+s+t=n} (-1)^{rs+t} m_{r+1+t}(\mathrm{id}^{\otimes r} \otimes m_s
  \otimes \mathrm{id}^{\otimes t}) = 0
  \qquad \forall\, n \geq 1.
\]

\begin{proposition}
  The commutator $[\alpha, \beta] = m_2(\alpha,\beta) - m_2(\beta,\alpha)
  \in \mathcal{A}^5$ is a boundary iff $m_3(\alpha,\alpha,\beta) = 0$ in
  $H^{-1}$ of the de Rham complex.  More generally, the vanishing of
  $m_n(\alpha^{n-1},\beta)$ is equivalent to a class in $H^{2-n}$ being zero.
\end{proposition}

\begin{proof}
  This is the standard obstruction theory for $A_\infty$-deformations, see
  Keller \cite{Keller2001}.  The class of $m_n$ lives in
  $\HH^3(\mathcal{A},\mathcal{A})$ (Hochschild cohomology), which in the
  bar-resolution picture corresponds to the degree $2-n$ part of $H^k$.
\end{proof}

\begin{theorem}\label{thm:infinite}
  Since $\dim H^k \geq 142$ for all $k$ in the computed truncation window
  $k \in [-2, 4]$, none of the cohomology groups vanish.  Therefore:
  \begin{enumerate}
    \item[\rm(i)] $m_n$ is potentially non-trivial for \emph{all} $n \geq 3$.
    \item[\rm(ii)] The $A_\infty$-deformation space of $\mathcal{A}$ is
          infinite-dimensional.
    \item[\rm(iii)] There is no finite presentation of $\mathcal{A}$ as a
          strict dg-algebra or as a path algebra of any quiver with relations.
  \end{enumerate}
\end{theorem}

\subsection{The self-referential obstruction}

Theorem~\ref{thm:infinite} reveals the self-referential structure that
makes $H^k$ an open problem:

\medskip
\begin{center}
\framebox{\parbox{0.85\textwidth}{\centering
  To classify $\gE$-modules: need $H^k$. \\
  $H^k$ encodes the $A_\infty$-maps $m_n$. \\
  The $m_n$ ARE the classification data. \\
  The problem is its own input.
}}
\end{center}
\medskip

This fixed-point structure is not circular in a logically vicious way.
It is, precisely, the structure of an \emph{undecidable} problem: no
finite computation resolves it because any resolution would encode all
resolutions of higher-order instances.

% =====================================================================
\section{Placement in the Descriptive Hierarchy}\label{sec:hierarchy}
% =====================================================================

\subsection{Review of the projective hierarchy}

The \emph{projective hierarchy} stratifies sets of real numbers (or more
generally of sequences of reals) by the complexity of their definitions:
\begin{align*}
  \Sigma_1^0 &= \text{open sets}, \quad \Pi_1^0 = \text{closed sets},\\
  \Sigma_{n+1}^1 &= \text{projections of $\Pi_n^1$ sets},\\
  \Pi_n^1 &= \text{complements of $\Sigma_n^1$ sets}.
\end{align*}
The union $\bigcup_n \Sigma_n^1$ is the \emph{projective sets}, also called
the \emph{analytic} sets in the broad sense.

The key levels for our purposes:
\begin{itemize}
  \item $\Sigma_1$ (arithmetic): the halting problem; quantifiers over $\NN$.
  \item $\Sigma_1^1$ (analytic): Borel-complete; one existential quantifier
        over $\NN^\NN$.  Classifying pairs of matrices under simultaneous
        similarity lands here (Belitskii--Sergeichuk \cite{BS2000}).
  \item $\Pi_1^1$ (co-analytic): the complement of $\Sigma_1^1$; universal
        statements over $\NN^\NN$.  Deciding $A_\infty$ quasi-isomorphism
        ($\forall n: m_n = 0$) lands here.
  \item Cofinal: not at any fixed level; reduces every analytic equivalence
        relation.
\end{itemize}

\subsection{Borel reducibility}

An equivalence relation $E$ on a Polish space $X$ is \emph{Borel-reducible}
to an equivalence relation $F$ on $Y$ (written $E \leq_B F$) if there
exists a Borel function $f: X \to Y$ with $x \mathrel{E} y \iff f(x) \mathrel{F} f(y)$.
The Borel reducibility order $(\leq_B, =_B)$ stratifies the complexity of
classification problems.

Key theorems:
\begin{itemize}
  \item Kechris--Louveau \cite{KL1997}: there is no maximum element in the
        Borel reducibility order (the hierarchy is cofinal).
  \item Hjorth \cite{Hjorth2000}: orbit equivalence relations of wild
        algebraic structures are not classifiable by countable structures.
  \item Louveau--Velickovic \cite{LV1994}: for every analytic $E$, there
        exists $E'$ with $E <_B E'$ (strict increase in complexity).
\end{itemize}

\subsection{Placement of NS and E(4,4)}

\begin{proposition}
  The following equivalence relations are ordered by strict Borel reducibility:
  \[
    E_{\mathrm{halt}} <_B E_{\mathrm{wild}} <_B E_{A_\infty} <_B E_{\gE}.
  \]
  Here:
  \begin{itemize}
    \item $E_{\mathrm{halt}}$: halting equivalence (Turing machine, same halting
          behavior); $\Sigma_1$-complete.
    \item $E_{\mathrm{wild}}$: simultaneous matrix similarity for all $n$;
          $\Sigma_1^1$-complete (Belitskii--Sergeichuk \cite{BS2000}).
    \item $E_{A_\infty}$: quasi-isomorphism of infinite $A_\infty$-algebras;
          $\Pi_1^1$-complete.
    \item $E_{\gE}$: isomorphism of finitely-generated $\gE$-modules;
          cofinal.
  \end{itemize}
\end{proposition}

% =====================================================================
\section{Cofinality of the NS Classification Problem}\label{sec:cofinality}
% =====================================================================

\begin{theorem}[Cofinality]\label{thm:cofinal}
  For every analytic equivalence relation $E$ on a Polish space, there
  is a Borel reduction from $E$ to the NS-solution equivalence relation
  $E_{\mathrm{NS}}$ (two initial data are equivalent if they produce
  NS-solutions in the same $\gE$-module isomorphism class).
\end{theorem}

\begin{proof}[Proof sketch]
  By Theorems~\ref{thm:wild} and~\ref{thm:infinite}, the $\gE$-module
  classification problem contains both Borel-complete problems ($\Sigma_1^1$,
  via the wild sub-quiver $Q_0$) and co-analytic problems ($\Pi_1^1$, via
  $A_\infty$ quasi-isomorphism).

  The key point is that the $A_\infty$-deformation space is
  \emph{parameterized by a Polish space}: the space of all choices of
  $\{m_n\}_{n \geq 3}$ satisfying the Stasheff relations is a closed subset
  of $\prod_{n \geq 3} \mathrm{Hom}(\mathcal{A}^{\otimes n}, \mathcal{A})$,
  which carries the product topology of a Polish space.

  By the Louveau--Velickovic theorem \cite{LV1994}, applied iteratively:
  for each $n$, the equivalence relation $E_n$ of isomorphism of
  $A_\infty$-algebras with $m_k = 0$ for all $k > n$ is analytic, and
  $E_n <_B E_{n+1}$ strictly.  The full $E_\gE = \bigcup_n E_n$ (modulo
  a density argument via Hjorth turbulence \cite{Hjorth2000}) is cofinal.

  The map $E \to E_{\mathrm{NS}}$ is realized by encoding the equivalence
  relation $E$ into the initial-data space $H^s(\RR^4)$ via the three
  surviving seed modes (Theorem~\ref{thm:main} of \cite{Paper1}), using
  the fact that the seed sector carries a faithful representation of
  $\mathcal{A}$ (Theorem~\ref{thm:obstruction}).
\end{proof}

\begin{corollary}
  No consistent extension of ZFC by countably many axioms can provide a
  complete system of Borel-measurable invariants for NS solutions.
\end{corollary}

\begin{proof}
  Any complete invariant system for $E_\mathrm{NS}$ would provide an upper
  bound on $E_\mathrm{NS}$ in the Borel reducibility order.  By cofinality,
  no such upper bound exists.
\end{proof}

% =====================================================================
\section{The Specific Arrow: $\phi_{4H}$}\label{sec:phi4H}
% =====================================================================

We isolate the single morphism responsible for crossing from ordinary wild
into $A_\infty$-wild complexity.

\begin{theorem}[Isolation of $\phi_{4H}$]\label{thm:isolation}
  \begin{enumerate}
    \item[\rm(i)] \emph{Without $\phi_{4H}$}: Complex B reduced to morphisms
          of degrees 1 and 3 only.  The Ext algebra $\mathcal{A}' = \Ext^*$
          restricted to degree-$\leq 3$ generators is a path algebra of a
          finite quiver.  Drozd's tame-wild dichotomy applies; the
          classification problem is at most $\Sigma_1^1$.
    \item[\rm(ii)] \emph{With $\phi_{4H}$}: The degree-4 generator $\beta$
          is added.  The Ext algebra $\mathcal{A}$ acquires generators in
          degrees 1 and 4.  The gcd $(1,4) = 1$, so the algebra is not
          Koszul.  The $A_\infty$-structure is unavoidable (Keller
          \cite{Keller2001}: any non-formal dg-algebra with generators
          in coprime degrees carries non-trivial $m_n$ for all $n$).
    \item[\rm(iii)] The degrees $1$ and $4$ are not arbitrary: they are
          determined by the singular-vector theory of $\gE$.
          Proposition 5.5 of CCK 2026 \cite{CCK2026} proves that singular
          vectors exist only in degrees 1, 2, 3, 4 for $\gE$.  The
          degree-4 morphism $\phi_{4H}$ is the maximum; it cannot be
          replaced by a lower-degree morphism without contradicting the
          singular-vector classification.
  \end{enumerate}
\end{theorem}

\begin{remark}
  All previously classified exceptional Lie superalgebras ($E(1|6)$,
  $E(3|6)$, $E(3|8)$) have exceptional de Rham complexes with morphisms
  only in degrees 1 and 2.  Their Ext algebras are Koszul, their
  $A_\infty$-deformation spaces are finite-dimensional, and their
  classification problems are $\Sigma_1^1$ (wild) but not $A_\infty$-wild.
  $\gE = E(4,4)$ is the first and (by the classification \cite{CCK2026})
  only exceptional Lie superalgebra for which a degree-4 morphism appears.
  It is, in this precise sense, genuinely exceptional even among
  exceptional algebras.
\end{remark}

% =====================================================================
\section{What Cannot Be Proved and Why}\label{sec:cannot}
% =====================================================================

The following statements are \emph{not merely open}: they are provably
unprovable within any fixed formal system $F$ that is consistent with ZFC.

\begin{theorem}\label{thm:unprovable}
  Let $F$ be any consistent formal system extending ZFC.  Then:
  \begin{enumerate}
    \item[\rm(i)] $F$ cannot prove ``there exists a complete system of
          Borel invariants for NS solutions,'' because such a theorem
          contradicts cofinality (Theorem~\ref{thm:cofinal}).
    \item[\rm(ii)] $F$ cannot prove ``$H^k(E(4,4)\text{-de Rham}) = V_k$
          for all $k$ simultaneously'' for any specific family of spaces
          $V_k$, because this would resolve the $A_\infty$-classification,
          which is cofinal.
    \item[\rm(iii)] $F$ cannot prove ``NS regularity is equivalent to
          [specific property $P$ of initial data]'' for any $P$ at a fixed
          level $\Sigma_n^1$, because NS regularity Borel-reduces to
          something harder than $\Sigma_n^1$ for every $n$.
  \end{enumerate}
\end{theorem}

\begin{remark}[The meta-level]
  The argument does not apply to Theorem~\ref{thm:unprovable} itself,
  because that theorem is a $\Pi_1^1$ statement (it is a universal claim
  about all formal systems $F$).  The statement ``no $F$ can prove X'' is
  of higher complexity than ``X''.  This is consistent with Gödel: each
  level of the hierarchy can prove statements about lower levels.
\end{remark}

% =====================================================================
\section{The New Complexity Stratum: $A_\infty$-Wild}\label{sec:newstratum}
% =====================================================================

We propose a name and formal definition for the new complexity level.

\begin{definition}[$A_\infty$-wild]\label{def:awild}
  A classification problem $\mathcal{C}$ (an equivalence relation on a
  Polish space of mathematical objects) is \emph{$A_\infty$-wild} if:
  \begin{enumerate}
    \item[\rm(i)] $\mathcal{C}$ is $\Sigma_1^1$-hard (Borel-complete);
    \item[\rm(ii)] $\mathcal{C}$ is $\Pi_1^1$-hard (co-analytic complete);
    \item[\rm(iii)] The classification data for $\mathcal{C}$ is naturally
          organized as an $A_\infty$-algebra (or $A_\infty$-category) with
          infinitely many non-trivial higher products $m_n$;
    \item[\rm(iv)] $\mathcal{C}$ is cofinal in the Borel reducibility order.
  \end{enumerate}
\end{definition}

\begin{theorem}
  The NS-solution classification problem via $\gE$ is $A_\infty$-wild.
\end{theorem}

\begin{proposition}[Hierarchy]
  $A_\infty$-wild is a strictly higher stratum than Borel-complete wild
  ($\Sigma_1^1$), which is strictly higher than the halting problem ($\Sigma_1$):
  \[
    \Sigma_1 <_B \Sigma_1^1\text{-wild} <_B \Pi_1^1\text{-wild}
    <_B A_\infty\text{-wild}.
  \]
\end{proposition}

The separations are witnessed by:
\begin{itemize}
  \item $\Sigma_1 <_B \Sigma_1^1$: well-known (any $\Sigma_1^1$-complete
        set is not $\Sigma_1$).
  \item $\Sigma_1^1 <_B \Pi_1^1$: $\Pi_1^1$ is not $\Sigma_1^1$-measurable
        (definable using a $\Sigma_1^1$ formula) in general.
  \item $\Pi_1^1 <_B A_\infty$-wild: cofinality (Theorem~\ref{thm:cofinal})
        implies that $A_\infty$-wild cannot be bounded by $\Pi_1^1$.
\end{itemize}

% =====================================================================
\section{Discussion}\label{sec:discussion}
% =====================================================================

\subsection{The role of $\phi_{4H}$}

The degree-4 exceptional morphism $\phi_{4H}$ is what elevates the problem
from ordinary wild to $A_\infty$-wild.  Its existence is forced by the
singular-vector theory: CCK 2026 proves that $\phi_{4H}$ corresponds to a
singular vector in $M_{t-4}(1,0,0)$ annihilated by all of $\mathfrak{g}_{+1}$,
with explicit formula involving fourth-order differential operators.  This
morphism cannot be removed or reduced in degree without contradicting the
algebraic structure of $\gE$.

\subsection{Open problems}

\begin{enumerate}
  \item \emph{Explicit $m_3$}: compute the leading $A_\infty$-product
        $m_3(\alpha,\alpha,\beta)$ from the de Rham complex differential.
        This requires determining $H^{-1}$ of the complex explicitly.
  \item \emph{Formal dg-algebra structure}: Is $\mathcal{A}$ formal
        (quasi-isomorphic to its cohomology as a dg-algebra)?  Formality
        would reduce the $A_\infty$-module classification to a single
        level $\Sigma_1^1$, contradicting cofinality unless $H^k = 0$
        for large $|k|$.  This is expected to fail.
  \item \emph{Large cardinal axioms}: $A_\infty$-wild problems may require
        large cardinal axioms (e.g., projective determinacy) to resolve
        individual instances.  Clarifying which large cardinal axiom level
        is needed is an open problem in mathematical logic.
\end{enumerate}

% =====================================================================
\section{The Navier--Stokes Liar's Paradox}\label{sec:liar}
% =====================================================================

The preceding sections show that $A_\infty$-wildness places NS above
every fixed level of the projective hierarchy.  We now show the logical
structure is sharper: the assumption of global smoothness is not merely
unprovable --- it is \emph{self-defeating} in precisely the sense of the
classical liar's paradox.

\subsection{The liar sentence}

Let $S$ denote the following sentence about NS:
\[
  S \;=\; \text{``NS does not simulate any formal system that decides $S$.''}
\]

\begin{theorem}[NS Liar's Paradox]\label{thm:liar}
  Assume ZFC is consistent.  Then the following are equivalent:
  \begin{enumerate}
    \item[\rm(a)] 4D NS has a global smooth solution for all smooth initial data;
    \item[\rm(b)] $\mathcal{A} = \Ext^*(M_*(1,0,0), M_*(1,0,0))$ is a formal dg-algebra;
    \item[\rm(c)] The NAND gate margin $\Delta$ (\cite[Thm.~2]{Paper1}) equals zero;
    \item[\rm(d)] NS does not simulate NAND gates;
    \item[\rm(e)] The Picard series $\sum_{n \geq 1} t^{n-1} C_n(k)$ converges for all $t \geq 0$.
  \end{enumerate}
  Each of \emph{(a)--(e)} implies every other.
  We have proved $\neg$(d) constructively: $\Delta = 0.007588 > 0$
  by interval arithmetic (\texttt{mpmath.iv}, $60$ significant digits).
  Therefore $\neg$(a) $\wedge$ $\neg$(b) $\wedge$ $\neg$(c) $\wedge$ $\neg$(e).
\end{theorem}

\begin{proof}[Proof sketch]
The implications form a chain:
\[
\text{(a)} \Rightarrow \text{(e)} \Rightarrow \text{(b)} \Rightarrow
\text{(b$'$: } \dim H^k \text{ bounded)} \Rightarrow
\text{classifiable at some } \Sigma_1^n \Rightarrow
\exists \text{ decision procedure } F.
\]
Any such $F$ is a Boolean circuit, hence a composition of NAND gates.
By \cite[Thm.~2]{Paper1} (NAND circuit simulation),
$F$ is embedded in the NS dynamics: NS simulates $F$.
But $F$ decides a statement about NS smoothness, so NS contains a
formal proof of its own global regularity.
By G\"odel's Second Incompleteness Theorem, the ambient formal system is
inconsistent, contradicting our assumption that ZFC is consistent.
Therefore (a) is false.  The equivalences then propagate via the
Blowup-Formality Theorem (Theorem~\ref{thm:blowup}).
\end{proof}

\subsection{Comparison with classical self-reference theorems}

\begin{table}[h]
\centering
\begin{tabular}{llll}
\toprule
\textbf{Theorem} & \textbf{Sentence} & \textbf{Resolution} & \textbf{Type} \\
\midrule
G\"odel (1931) & ``This is not provable in $F$.'' &
  True but unprovable & Incompleteness \\
Tarski (1936) & ``This is not true in $M$.'' &
  Truth not definable in $M$ & Undefinability \\
\textbf{NS Liar} & $S$ above &
  Smoothness is \textbf{refuted} & \textbf{Refutation} \\
\bottomrule
\end{tabular}
\caption{Classical self-reference theorems compared.  The NS Liar's Paradox
  is the unique case where the self-referential sentence yields a
  constructive disproof rather than unprovability or undefinability.}
\label{tab:liar}
\end{table}

\medskip

The key distinction: G\"odel yields \emph{unprovability}; the NS Liar yields
a \emph{disproof}.  We do not merely lack a proof of smoothness.
The interval certificate $\Delta = 0.007588 > 0$ is a ZFC theorem that
concretely refutes the assumption leading to smoothness.

\begin{remark}
  The NAND initial conditions of \cite[Thm.~2]{Paper1} are
  Lawvere's fixed point made concrete: they simultaneously
  (i)~perform a computation (evaluate NAND), (ii)~witness Turing universality,
  and (iii)~generate, via the chain above, evidence of their own finite-time
  blowup.  The blowup at $T_{\mathrm{crit}}$ is the liar's paradox resolving.
\end{remark}

\begin{remark}[The general principle]
  No system of logic can make truth statements about self-similar
  singularities or blowups.  The NS Liar's Paradox is one instance;
  the self-similar fixed point of the renormalization-group flow for
  a QCD glueball is another; the fixed point of a self-improving
  intelligence acting on its own architecture is a third.  In each case
  the object is the fixed point of its own description, and Lawvere's
  fixed-point theorem (the abstract form of the liar) guarantees that
  no consistent external logic can certify it.
\end{remark}

% =====================================================================
\begin{thebibliography}{99}
% =====================================================================

\bibitem{CCK2026}
N.~Cantarini, F.~Caselli, and V.~Kac,
\emph{Classification of degenerate Verma modules over $E(4,4)$},
arXiv:2603.16507, 2026.

\bibitem{Kac1982}
V.~G.~Kac,
\emph{Root systems, representations of quivers and invariant theory},
in \emph{Invariant Theory}, Lecture Notes in Mathematics \textbf{996},
Springer, 1983, 74--108.

\bibitem{Gabriel1972}
P.~Gabriel,
\emph{Unzerlegbare Darstellungen I},
Manuscripta Math.\ \textbf{6} (1972), 71--103.

\bibitem{Drozd1979}
Ju.~A.~Drozd,
\emph{Tame and wild matrix problems},
in \emph{Representations and Quadratic Forms}, Kiev, 1979, 39--74.

\bibitem{BS2000}
G.~R.~Belitskii and V.~V.~Sergeichuk,
\emph{Complexity of matrix problems},
Linear Algebra Appl.\ \textbf{361} (2003), 203--222.

\bibitem{Stasheff1963}
J.~D.~Stasheff,
\emph{Homotopy associativity of $H$-spaces I, II},
Trans.\ Amer.\ Math.\ Soc.\ \textbf{108} (1963), 275--312.

\bibitem{Keller2001}
B.~Keller,
\emph{Introduction to $A$-infinity algebras and modules},
Homology Homotopy Appl.\ \textbf{3} (2001), 1--35.

\bibitem{Hjorth2000}
G.~Hjorth,
\emph{Classification and Orbit Equivalence Relations},
Mathematical Surveys and Monographs \textbf{75},
American Mathematical Society, 2000.

\bibitem{KL1997}
A.~S.~Kechris and A.~Louveau,
\emph{The classification of hypersmooth Borel equivalence relations},
J.\ Amer.\ Math.\ Soc.\ \textbf{10} (1997), 215--242.

\bibitem{LV1994}
A.~Louveau and B.~Velickovic,
\emph{A note on Borel equivalence relations},
Proc.\ Amer.\ Math.\ Soc.\ \textbf{120} (1994), 255--259.

\bibitem{Bondal1990}
A.~I.~Bondal and M.~M.~Kapranov,
\emph{Enhanced triangulated categories},
Math.\ USSR Sb.\ \textbf{70} (1991), 93--107.

\bibitem{Paper1}
D.~Remmenga,
\emph{Reduction of the Four-Dimensional Navier--Stokes and Euler Problems
to Three Surviving Modes in $E(4,4)$},
companion paper, 2026.

\end{thebibliography}

\end{document}
